\chapter[Coincidence Classification]{Coincidences amongst families of mesh patterns and classical patterns}\label{chap:coincs}
One interesting question to ask about permutation patterns is to consider when a
pattern may be avoided by, or contained in, arbitrary permutations. Two
patterns \(\pi\) and \(\sigma\) are said to be \emph{coincident} if the set
of permutations that avoid \(\pi\) is the same as the set of permutations that
avoid \(\sigma\), \ie \(\av{\pi} = \av{\sigma}\). This extends to sets of patterns as well as single patterns.

We consider the avoidance sets, \(\av{\{\pi,q\}}\) where \(\pi\) is a classical pattern
of length \(3\) and \(q\) is a mesh pattern of length \(2\) in order to
establish some rules about when these two sets give the same avoidance set.
We fix \(\pi\) in order to define the  coincidence and say that \(\pi\) is the
\emph{dominating pattern}.

It is useful to be able to modify a mesh pattern by adding points to an
already existing mesh pattern. First adding a single point into a pattern.
\begin{definition}
Given a mesh pattern \(p = (\pi,R)\) \(\addp[D]{p}{(a,b)}\) gives a mesh
pattern \(p^\prime = (\pi^\prime,R^\prime)\) with length \(\setsize{\pi}+1\) defined by
\begin{equation*}
    \pi^\prime(i) = \begin{cases}
    \pi(i) & \text{if } i\neq a+1 \text{ and } \pi(i)<b\\
    \pi(i)+1 & \text{if } i\neq a+1 \text{ and } \pi(i)>b\\
    b+1 & \text{if } i = a+1
\end{cases}
\end{equation*}
and
\begin{equation*}
        R^\prime = \bigcup_{(i,j)\in R}r((i,j)),
\end{equation*}
where \(r((i,j))\) is defined by
\begin{equation*}
    r((i,j)) =
    \begin{cases}
        \{(i,j)\} & \text{if } i<a, j<b\\
        \{(i,j),(i,j+1)\} & \text{if } i<a, j=b\\
        \{(i,j+1)\} & \text{if } i<a, j>b\\
        \{(i,j),(i+1,j)\} & \text{if } i=a, j<b\\
        \{(i,j+1),(i+1,j+1)\} & \text{if } i=a, j>b\\
        \{(i+1,j)\} & \text{if } i>a, j<b\\
        \{(i+1,j),(i+1,j+1)\} & \text{if } i>a, j=b\\
        \{(i+1,j+1)\} &\text{if } i>a, j>b\\
    \end{cases}
\end{equation*}

If the shading set \(D\) is non-empty we can modify definition of the directions
slightly
\begin{equation*}
    \begin{aligned}
        N &= \{(a,b+1),(a+1,b+1)\} \\
        E &= \{(a+1,b),(a+1,b+1)\} \\
        S &= \{(a,b),(a+1,b)\} \\
        W &= \{(a,b),(a,b+1)\}
    \end{aligned}
\end{equation*}
And we add the union of the sets in \(D\) into the mesh set \(R^\prime\).
\end{definition}

Given a mesh pattern \(p\), \(\addp[D]{p}{(a,b)}\) is the operation that returns
a mesh pattern equivalent to placing a point in the center of box \((a,b)\),
where \((a,b)\) is not in the mesh set of \(p\) , with shading defined by
\(D\subseteq\{N,E,S,W\}\).
The set \(D\) defines the shading by indicating that the boxes in the cardinal
directions in \(D\) next to the point are shaded in the resulting pattern.
Since there is no ambiguity we let \(\addpe[D]\) be equivalent to \(\addp[D]{\varepsilon}{(0,0)}\).

\begin{example}
    The result of adding a single point to the empty permutation for each cardinal direction.
 \begin{align*}
     \addpe[\{N\}] = \textpattern{}{1}{0/1,1/1} &\qquad\addpe[\{E\}] =\textpattern{}{1}{1/1,1/0} \qquad \\
     \addpe[\{S\}] = \textpattern{}{1}{0/0,1/0} &\qquad\addpe[\{W\}] =\textpattern{}{1}{0/0,0/1} \qquad
\end{align*}
    A more complex example for \texttt{add\_point} could be
\begin{equation*}
\addp[\{N,E\}]{\pattern{}{2,1,3}{0/3}}{(2,3)} = \pattern{}{2,1,4,3}{0/3,0/4,2/4,3/4,3/3}
\end{equation*}
\end{example}

It is also useful to think about adding an ascent, or descent, into a pattern

\begin{definition}
Considering
only adding the ascent, as adding a descent is very similar.
Given a mesh pattern \(p = (\pi, R)\) \(\adda[D]{p}{(a,b)}\) gives a mesh
pattern \(p^\prime = (\pi^\prime,R^\prime)\) with length \(\setsize{\pi}+2\) defined by
\begin{equation*}
    \pi^\prime(i) = \begin{cases}
    \pi(i) & \text{if } i\neq a+1,a+2 \text{ and } \pi(i)<b\\
    \pi(i)+2 & \text{if } i\neq a+1,a+2 \text{ and } \pi(i)>b\\
    b+j, b+j & \text{if } i = a+j, j\in\{1,2\}
\end{cases}
\end{equation*}
and
\begin{equation*}
        R^\prime = \{(a+1,b),(a,b),(a+1,b),(a+2,b),(a+1,b+2)\} \cup \bigcup_{(i,j)\in R}r((i,j)),
\end{equation*}
where \(r((i,j))\) is defined by
\begin{equation*}
    r((i,j)) =
    \begin{cases}
        \{(i,j)\} & \text{if } i<a, j<b\\
        \{(i,j),(i,j+1),(i,j+2)\} & \text{if } i<a, j=b\\
        \{(i,j+2)\} & \text{if } i<a, j>b\\
        \{(i,j),(i+1,j),(i+2,j)\} & \text{if } i=a, j<b\\
        \{(i,j+2),(i+1,j+2),(i+2,j+2)\} & \text{if } i=a, j>b\\
        \{(i+2,j)\} & \text{if } i>a, j<b\\
        \{(i+2,j),(i+2,j+1),(i+2,j+2)\} & \text{if } i>a, j=b\\
        \{(i+2,j+2)\} &\text{if } i>a, j>b\\
    \end{cases}
\end{equation*}
\end{definition}

Given a pattern \(p\), \(\addd{p}{(a,b)}\), and \\ \(\adda{p}{(a,b)}\),
are the operations that return a mesh pattern equivalent to placing a
decrease, or increase, in the center of box \((a,b)\), where \((a,b)\) is not
in the mesh set of \(p\), in \(p\).

\begin{example}
    \begin{align*}
        \addae = \pattern{}{1,2}{0/1,1/0,1/1,1/2,2/1}\\
        \addde = \pattern{}{2,1}{0/1,1/0,1/1,1/2,2/1}
    \end{align*}

    \noindent A more complex example is
\begin{equation*}
\adda{\pattern{}{2,3,1}{0/0,1/0,1/3,
                        3/0,3/1,3/3}}{(1,1)}
= \pattern{}{4,2,3,5,1}{0/0,
                        1/0,1/2,1/5,
                        2/0,2/1,2/2,2/3,2/5,
                        3/0,3/2,3/5,
                        5/0,5/1,5/2,5/3,5/5}
\end{equation*}
\end{example}

We now attempt to fully classify coincidences in families characterised by avoidance
of a classical pattern of length \(3\) and a mesh pattern of length \(2\), that is
finding and explaining all coincidences where \(\av{\{p,m\}} = \av{\{p,m^\prime\}}\).

It can be easily seen that in order to classify coincidences one need only
consider coincidences within the family of mesh patterns with the same underlying
classical pattern, this is due to the fact that \(\perm{2,1}\in\av{\mperm{1,2}{R}}\)
and \(\perm{1,2}\in\av{\mperm{2,1}{R}}\) for all mesh-sets \(R\).

We know that there are a total of \(512\) mesh-sets for each underlying classical
pattern. By use of the previous results of \textcite{DBLP:journals/corr/ClaessonTU14}\footnote{
The authors use the Simultaneous Shading Lemma, a closure result and one worked out special case.
} the number of coincidence classes can be reduced to \(220\).
% the Shading
% Lemma\cite[Lemma~3.11]{DBLP:journals/combinatorics/HilmarssonJSVU15},
% Simultaneous Shading Lemma\textcite[Lemma~7.6]{DBLP:journals/corr/ClaessonTU14}, and
% one special case, we can reduce the number of equivalence classes to \(220\).

By discussion of a number of rules we will show that the number of coincidence
classes follows the values shown in  \TableRef{tab:domclasses}. The experimental
data in the last row of the table is calculated on permutations up to length 11.

\begin{table}[htb]
\begin{center}
\begin{tabularx}{\textwidth}{c|Y|Y|Y|Y|}
\cline{2-5}
& \multicolumn{4}{|c|}{Dominating Pattern}\\
\cline{2-5}
& \multicolumn{2}{|c|}{\(\perm{2,3,1}\)}&\multicolumn{2}{|c|}{\(\perm{3,2,1}\)}\\
\cline{2-5}
& \(\perm{1,2}\)&\(\perm{2,1}\)&\(\perm{1,2}\)&\(\perm{2,1}\)\\
\hline
\multicolumn{1}{|r|}{No Dominating rule}&\(220\)&\(220\)&\(220\)&\(220\)\\
\hline
\multicolumn{1}{|r|}{First Dominating rule}&\(85\)&\(43\)&\(220\)&\(29\)\\
\hline
\multicolumn{1}{|r|}{Second Dominating rule}&\(59\)&\(39\)&\(220\)&\(29\)\\
\hline
\multicolumn{1}{|r|}{Third Dominating rule}&\(56\)&\(39\)&\(220\)&\(29\)\\
\hline
\multicolumn{1}{|r|}{Experimental class size}&\(56\)&\(39\)&\(213\)&\(29\)\\
\hline
\end{tabularx}
\end{center}
    \caption{Coincidence class number reduction by application of Dominating Rules}
    \label{tab:domclasses}
\end{table}

From the table it can be seen that the rules established capture almost all
coincidences. However, there are still some coincidences that are not able to
be explained by the rules. This shows that complete coincidence classification
of mesh patterns is a very difficult task, even when we have additional tools
available.
\section{Coincidence classes of Av(\{321, (21, \textit{R})\}).}
Through experimentation, up to permutations of length 11, we discover that there
are at least 29 coincidence classes of mesh patterns with underlying classical
pattern \(\perm{2,1}\).

\begin{proposition}[First Dominating Pattern Rule]
    \label{prop:dom1}
    Given two mesh patterns \(m_1 =(\sigma, R_1)\) and \(m_2 = (\sigma, R_2)\),
    and a dominating classical pattern \(\pi = (\pi,\emptyset)\) such that
    \(\setsize{\pi} \le \setsize{\sigma} + 1\), the sets \(\av{\{\pi,m_1\}}\) and
    \(\av{\{\pi,m_2\}}\) are coincident if

    \begin{enumerate}
        \item \(R_1 \triangle R_2 = \{(a,b)\}\)
        \item \(\pi \preceq \addp{\sigma}{(a,b)}\)\label{prop:dom1:cont}
    \end{enumerate}
\end{proposition}
In order to prove this proposition we must first make the following note.

\begin{note}
    \label{not:downcmesh}
    Let \(R^\prime \subseteq R\). Then any occurence of \((\tau, R)\) in a permutation
    is an occurence of \((\tau,R^\prime)\).
\end{note}
% \begin{proof}
%     Assume without meaningful loss of generality that \(R^\prime\) is a proper
%     subset of \(R\).
%
%     Consider an occurence of \((\tau, R)\) in a permutation \(\sigma\), obviously
%     this corresponds to an occurrence of \(\tau\) in \(\sigma\). Now consider
%     the mesh sets \(R\) and \(R^\prime\), since \(R^\prime \subseteq R\) then
%     there are more restrictions on where points are in an occurrence of \((\tau,R)\).
%     Namely, for every shaded box in \(R\) the corresponding region in \(\sigma\)
%     must contain no points, since \(R^\prime\) has less shading than \(R\) there
%     exists a region in the occurence of \((\tau,R)\) in \(\sigma\) that is
%     now devoid of restrictions. However, by removing restrictions we cannot
%     make an occurence become not an occurence, and therefore the same occurrence
%     of \(\tau \) in \(\sigma\) is now an occurence of \((\tau, R^\prime)\).
% \end{proof}
\begin{proof}[Proof of \PropositionRef{prop:dom1}]
    We need to prove that \(\av{\{\pi,m_1\}} = \av{\{\pi,m_2\}}\).

    \noindent Assume without meaningful loss of generality that \(R_2 = R_1 \cup \{(a,b)\}\).
    % Consider a permutation \(\omega\) that contains an occurrence of \(m_2\)
    % by \NoteRef{not:downcmesh} any of these occurrences is also an occurence
    % of \(m_1\). This proves that every occurence of \(m_2\) is also an
    % occurence of \(m_1\) and therefore \(\av{\{\pi,m_1\}} \subseteq \av{\{\pi,m_2\}}\).
    Since \(R_1\) is a subset of \(R_2\), \NoteRef{not:downcmesh} states that
    \(\av{\{\pi,m_1\}} \subseteq \av{\{\pi,m_2\}}\)

    Now we consider a permutation \(\omega^\prime\in \av{\pi}\),
    containing an occurrence of \(m_1\). Consider placing a point in the
    region corresponding to the box \((a,b)\), regardless of where in this
    region we place the point by condition~\ref{prop:dom1:cont} of the
    Proposition we create an occurrence of \(\pi\), therefore there can be no
    points in this region, which could have been represented in the mesh set
    \(R_1\) by adding the box \((a,b)\). Hence every occurrence of
    \(m_1\) is in fact an occurrence of \(m_2\), and we have that
    \(\av{\{\pi,m_2\}} \subseteq \av{\{\pi,m_1\}}\).

    Taking both directions of the containment we can therefore draw the
    conclusion that \(\av{\{\pi,m_1\}} = \av{\{\pi,m_2\}}\).
\end{proof}

All coincidence classes of \(\av{\{\perm{3,2,1}, \mperm{2,1}{R}\}}\) can be
explained by application of \PropositionRef{prop:dom1}. % repeated?

This rule can be understood very in graphical form. In the
pattern in \FigureRef{fig:rule1} we can gain shading in three boxes since if
there is a point in any of these boxes we would get an occurrence of the
dominating pattern \(\perm{3,2,1}\).

\begin{figure}[htb]
    \begin{center}
    \begin{tikzpicture}[scale=\picscale]
        \modpattern{}{2,1}{1/0,1/1,1/2,2/2}
    \end{tikzpicture}
        \raisebox{2ex}{\(\mapsto\)}
    \begin{tikzpicture}[scale=\picscale]
        \filldraw (2.5,0.5) circle (3pt);
        \modpattern{}{2,1}{1/0,1/1,1/2,2/2,2/0}
    \end{tikzpicture}

        \caption{Visual depiction of first dominating pattern rule.}
        \label{fig:rule1}
    \end{center}
\end{figure}

Note that there are two natural extensions of this rule. We can replace \(\pi\)
with a set of classical patterns, or we can consider \(\pi\) to be a mesh pattern.

\section{Equivalence classes of Av(\{231, (21, \textit{R})\}).}
By application of \PropositionRef{prop:dom1} we obtain \(43\) coincidence
classes. Experimentation shows that there are in fact at least \(39\) coincidence
classes, for example the following two patterns are coincident in
\(\av{\perm{2,3,1}}\) but this is not explained by \PropositionRef{prop:dom1}.
\begin{equation*}
    m_1 = \pattern{}{2,1}{0/0,0/2,1/0,2/1} \text{ and } m_2 = \pattern{}{2,1}{0/0,0/2,1/0,1/1,2/1}
\end{equation*}

Consider an occurrence of \(m_1\) in a permutation in \(\av{\perm{2,3,1}}\), consisting
of elements \(x\) and \(y\). If the region corresponding
to the box \((1,1)\) is empty we have an occurrence of \(m_2\). Otherwise, if there is any
increase in this box then we would have an occurrence of \(\perm{2,3,1}\),
however, since we are in \(\av{\perm{2,3,1}}\) this is not possible. This box
must therefore contain a (non-empty) decreasing subsequence.
 This gives rise to the following lemma:

\begin{lemma}
    \label{lem:incdec}
    Given a mesh pattern \(m =(\sigma, R)\), where the box \((a,b)\) is not
    in \(R\), and a dominating classical pattern \(\pi = (\pi,\emptyset)\) if
    \(\pi \preceq \adda{\sigma}{(a,b)}\)\\(\(\pi \preceq \addd{\sigma}{(a,b)}\))
    then in any occurrence of \(m\) in a permutation \(\varrho\) the region
    corresponding to the box \((a,b)\) can only contain an increasing
    (decreasing) subsequence of \(\varrho\).
\end{lemma}
The proof is analogous to the proof of \PropositionRef{prop:dom1}.

Going back to our example mesh patterns
\begin{equation*}
    \begin{tikzpicture}[scale=\picscale]
        \modpattern{}{2,1}{0/0,0/2,1/0,2/1}
        \filldraw (1.66,1.33) circle (3pt);
        \filldraw (1.33,1.66) circle (3pt);
    \end{tikzpicture}
\end{equation*}
We know that the region corresponding to the box \((1,1)\) contains a decreasing subsequence.
If we let \(z\) be the topmost point in this decereasing subsequence, then \(xz\)
is an occurrence of \(m_2\). This shows that our two example patterns are coincident.

This result generalises into the following rule for categorising coincidences
of mesh patterns in cases where there is a dominating classical pattern.

\begin{proposition}[Second Dominating Pattern Rule]
    \label{prop:dom2}
    Given two mesh patterns \(m_1 =(\sigma, R_1)\) and \(m_2 = (\sigma, R_2)\),
    and a dominating classical pattern \(\pi = (\pi,\emptyset)\) such that
    \(\setsize{\pi} \le \setsize{\sigma} + 2\), the sets \(\av{\{\pi,m_1\}}\) and
    \(\av{\{\pi,m_2\}}\) are coincident if

    \begin{enumerate}
        \item \(R_1 \triangle R_2 = \{(a,b)\}\)
        \item   \begin{enumerate}
                \item\label{prop:dom2:condc} \(\pi \preceq \adda{\sigma}{(a,b)}\) and
                        \begin{enumerate}
                            \item \((a+1,b) \in \sigma\) and \((a+1,b-1)\notin R\) and \\
                                \((x,b-1)\in R \implies (x,b) \in R \) (where \(x\neq a,a+1\)) and\\
                                  \((a+1,y)\in R \implies (a,y) \in R\) (where \(y\neq b-1,b\)).
                            \item \((a,b+1) \in \sigma\) and \((a-1,b+1)\notin R\) and \\
                                  \((x,b+1)\in R \implies (x,b) \in R\) (where \(x\neq a-1,a\)) and\\
                                  \((a-1,y)\in R \implies (a,y) \in R\) (where \(y\neq b,b+1\)).
                        \end{enumerate}
                    \item \(\pi \preceq \addd{\sigma}{(a,b)}\) and
                        \begin{enumerate}
                            \item \((a+1,b+1) \in \sigma\) and \((a+1,b+1)\notin R\) and \\
                                  \((x,b+1)\in R \implies (x,b) \in R\) (where \(x\neq a,a+1\)) and\\
                                  \((a+1,y)\in R \implies (a,y) \in R\) (where \(y\neq b,b+1\)).
                            \item \((a,b) \in \sigma\) and \((a-1,b-1)\notin R\) and \\
                                  \((x,b+1)\in R \implies (x,b) \in R\) (where \(x\neq a-1,a\)) and\\
                                  \((a-1,y)\in R \implies (a,y) \in R\)  (where \(y\neq b-1,b\)).
                        \end{enumerate}
                \end{enumerate}
    \end{enumerate}
\end{proposition}
\begin{proof}
    We need to prove that \(\av{\{\pi,m_1\}} = \av{\{\pi,m_2\}}\).

    \noindent Assume without meaningful loss of generality that \(R_2 = R_1 \cup \{(a,b)\}\).

    Consider a permutation \(\omega\) that contains an occurrence of \(m_2\).
    By \NoteRef{not:downcmesh} any of these occurrences is also an occurence
    of \(m_1\). This proves that every occurence of \(m_2\) is also an
    occurence of \(m_1\) and therefore \(\av{\{\pi,m_1\}} \subseteq \av{\{\pi,m_2\}}\).

    We will consider taking the first branch of every choice.
    Now consider a permutation \(\omega^\prime \in \av{\pi}\). Suppose
    \(\omega^\prime\) contains \(m_1\) and consider the region corresponding to
    \((a,b)\) in \(R_1\).

    If the region is empty, the occurrence of \(m_1\) is trivially an occurence
    of \(m_2\).

    Now consider if the region is non-empty, by \LemmaRef{lem:incdec} and
    condition~\ref{prop:dom2:condc} of the proposition this region must contain
    a decreasing subsequence.
    We can choose the topmost point in the region to replace the corresponding
    point in the mesh pattern and the points from the subsequence are now in
    the box southeast of the point. The other conditions allow this to be done
    without points being present in regions that were shaded. Hence there are no
    points in the region corresponding to the box \((a,b)\) in the mesh
    pattern, and therefore we can shade this region. This implies that every
    occurrence of \(m_1\) in \(\av{\pi}\) is in fact an occurrence of \(m_2\)
    so \(\av{\{\pi,m_2\}} \subseteq \av{\{\pi,m_1\}}\).

    Similar arguments satisfy the remainder of the branches.
\end{proof}
\begin{figure}[ht]
\centering
\begin{tikzpicture}[scale = 0.75]

     % Make the grid
     \draw[step=1cm,gray!50] (2/3,2/3) grid (5+1/3,5+1/3);

     % Make black cross in the middle
     \draw[-,very thick] (3,2/3) -- (3,5+1/3);
     \draw[-,very thick] (2/3,3) -- (5+1/3,3);
     \draw[-,very thick,gray] (4,2) -- (2,2) -- (2,4) -- (4,4) -- (4,2);

     \fill[black] (3,3) circle (0.1);
     \fill[red] (2,3) circle (0.1);

     % Draw the dots on the side
     \foreach \x in {1,2,4,5}{
         \draw[dotted, gray!50] (\x, 0) -- (\x, 2/3);
         \draw[dotted, gray!50] (\x, 5+1/3) -- (\x, 6);
         \draw[dotted, gray!50] (0,\x) -- (2/3,\x);
         \draw[dotted, gray!50] (5+1/3,\x) -- (6,\x);
        }
     % Draw black dots next to the cross
     \draw[dotted, black, very thick] (3, 0) -- (3, 2/3);
     \draw[dotted, black, very thick] (3, 5+1/3) -- (3, 6);
     \draw[dotted, black, very thick] (0, 3) -- (2/3, 3);
     \draw[dotted, black, very thick] (5+1/3, 3) -- (6, 3);

     \draw[->-=.5,thick] (3,3) to[bend right=15](2,4);

     \draw[->, thick] (5,3) -- (5,4);
     \draw[->, thick] (1,3) -- (1,4);
     \draw[->, thick] (3,5) -- (2,5);
     \draw[->, thick] (3,1) -- (2,1);

     \node [black,right] at (6,3) {$b$};
     \node [black,below] at (2,0) {$a$};

\end{tikzpicture}
    \caption{If the conditions of \PropositionRef{prop:dom2} are satisfied the box
    \((a-1,b)\) can be shaded.}
\end{figure}

This proposition essentially states that we slide all of the points in the box we desire
to shade diagonally, and chose the topmost/bottommost point to replace the original
point in the mesh pattern.

By taking the First Dominating Pattern Rule and the Second Dominating Pattern
Rule together coincidences of classes of the form \(\av{\{\perm{2,3,1},\mperm{2,1}{R}\}}\)
are fully explained, obtaining 39 coincidence classes of mesh patterns.

\section{Equivalence classes of Av(\{231, (12, \textit{R})\}).}
When considering the coincidence classes of \(\av{\perm{2,3,1},\mperm{1,2}{R}}\)
we first apply the two Dominating Pattern rules previously established.
Starting from \(220\) classes, application of the first Dominating Pattern rule
gives \(85\) classes. Following this with the second Dominating Pattern rule
reduces the number of classes to \(59\). However we know that there are patterns
where the coincidences are not explained by the rules given above.

For example the patterns
\begin{equation*}
    m_1 = \pattern{}{1,2}{0/2,2/0,2/1} \text{ and } m_2 = \pattern{}{1,2}{0/2,1/0,2/0,2/1}
\end{equation*}
are experimentally coincident. This coincidence is not explained by our rules,
it is necessary to attempt to capture these coincidences by establishing more
rules.

Consider an occurrence of \(m_1\) in a permutation, if the region corresponding
to the box \((1,0)\) is empty then we have an occurrence of \(m_2\). Now look at
the case when this region is not empty. Consider choosing the rightmost point in
region.
% In order to rigorously establish this coincidence we need to consider what would
% happen if we were to choose a point in the box \((1,0)\) that we would like to shade.
% In order to get a chance of getting the shading we want consider choosing the
% rightmost point in the box.
This gives us an occurrence of the following mesh pattern.
\begin{equation*}
    \pattern{}{2,1,3}{0/3,2/0,2/1,3/0,3/1,3/2}
\end{equation*}
By application of \PropositionRef{prop:dom1} we then achieve the following
mesh pattern
\begin{equation*}
    \pattern{2,3}{2,1,3}{0/3,1/3,2/0,2/1,3/0,3/1,3/2}
\end{equation*}
If we look at the highlighted points we see that the subpattern is an occurrence
of the mesh pattern that we originally desired. This gives rise to the following
rule:

\begin{proposition}[Third Dominating Pattern Rule]
    \label{prop:dom3}
    Given two mesh patterns \(m_1 =(\sigma, R_1)\) and \(m_2 = (\sigma, R_2)\),
    and a dominating classical pattern \(\pi = (\pi,\emptyset)\), the sets
    \(\av{\{\pi,m_1\}}\) and \(\av{\{\pi,m_2\}}\) are coincident if
    \begin{enumerate}
        \item \(R_1 \triangle R_2 = \{(a,b)\}\)
        \item\label{prop:dom3:condocc} \(\addp[D]{\mperm{\sigma}{R_1}}{(a,b)}\)
            where \(D\in\{N,E,S,W\}\)
            is coincident with a mesh pattern containing an occurrence of
            \(\mperm{\sigma}{R_2}\) as a subpattern. 
    \end{enumerate}
\end{proposition}
\begin{proof}
    We need to prove that \(\av{\{\pi,m_1\}} = \av{\{\pi,m_2\}}\).

    \noindent Assume without meaningful loss of generality that \(R_2 = R_1 \cup \{(a,b)\}\).

    Consider a permutation \(\omega\) that contains an occurrence of \(m_2\).
    By \NoteRef{not:downcmesh} any of these occurrences is also an occurence
    of \(m_1\). This proves that every occurence of \(m_2\) is also an
    occurence of \(m_1\) and therefore \(\av{\{\pi,m_1\}} \subseteq \av{\{\pi,m_2\}}\).

    Now consider a permutation \(\varrho\) in \(\av{\pi}\) that contains an
    occurrence of \(m_1\). If the region corresponding to the box \((a,b)\)
    is empty then we have an occurrence of \(m_2\). If the region is non-empty
    then by condition~\ref{prop:dom3:condocc}  of the proposition there exists
    a direction such that there exists an occurrence of a mesh pattern of
    length one longer than \(m_1\) in this position. This mesh pattern is
    coincident with another mesh pattern that contains an
    occurrence of \(m_2\). Hence, every occurrence of \(m_1\) leads to an occurrence
    of \(m_2\).  Thus \(\av{\{\pi,m_2\}} \subseteq \av{\{\pi,m_1\}}\) and the
    two patterns are coincident.
\end{proof}

By application of this rule we can reduce the number of classes in
\(\av{\{\perm{2,3,1},(\perm{1,2}, R)\}}\) to \(56\).

\section{Equivalence classes of Av(\{321, (12, \textit{R})\}).}
When considering coincidences of mesh patterns with underlying classical
pattern \(\perm{1,2}\) in \(\av{\perm{3,2,1}}\) application of the previously
established rules give no coincidences. Through experimentation we discover
that there are \(7\) non-trivial coincidence classes (all others are singletons)
which can be explained through the
use of two different lines of reasoning. Since the number of coincidences is
so small we will reason for these coincidences without attempting to generalise
into concrete rules.

Intuitively it is easy to see why our previous rules have no
power here. There is nowhere that it is possible to add a single point to
gain an occurrence of \(\pi = \perm{3,2,1}\). It is also impossible to have
a position where addition of an increase, or decrease, provides extra
shading power.

The patterns
\begin{equation*}
    m_1 = \pattern{}{1,2}{0/0, 0/2, 1/1, 1/2, 2/1} \text { and } m_2 \pattern{}{1,2}{0/0, 0/1, 0/2, 1/1, 1/2, 2/1}
\end{equation*}
are equivalent in \(\av{\perm{3,2,1}}\). (There are 3 symmetries of these
patterns that are also equivalent to each other by the same reasoning.)

Consider the region corresponding to the box \((0,1)\) in any occurrence of \(m_1\),
in a permutation. By \LemmaRef{lem:incdec} it must contain an increasing subsequence.
If the region is empty then we have an occurrence of \(m_2\).
If there is only one point in the region we can choose this to replace the \(1\) in
the mesh pattern to get the required shading. If there is more than one point then
choosing the two leftmost points gives us the following mesh pattern.
\begin{equation*}
    \pattern{3,4}{2,3,1,4}{0/0, 0/1, 0/2, 0/3, 0/4,
                            1/0, 1/1, 1/2, 1/3, 1/4,
                            2/0, 2/1,2/2, 2/4,
                            3/1, 3/2, 3/3, 3/4,
                            4/1, 4/2, 4/3}
\end{equation*}
Where the two highlighted points are the original two points.
Now if we take the other two points as the points in our
mesh pattern then we get an occurrence of the pattern
we originally desired, and hence the two patterns are coincident.

The other reasoning applies to the patterns
\begin{equation*}
    m_1 = \pattern{}{1,2}{0/1,0/2,2/0,2/1} \text{ and } m_2 = \pattern{}{1,2}{0/1,0/2,1/1,2/0,2/1}
\end{equation*}
which are coincident by experimentation.

In order to prove this coincidence we will proceed by mathematical
induction on the number of points in region corresponding to the middle box.
We call this number \(n\).
\begin{description}
    \item [Base Case \((n=0)\):] The base case holds since we can freely shade the
        box if it contains no points.
    \item [Inductive Hypothesis \((n=k)\):] Suppose that we can find an occurrence
        of the second pattern if we have an occurrence of the first with \(k\) points
        in the middle box.
    \item [Inductive Step \((n=k+1)\)] Suppose that we have \((k+1)\) points in
        the middle box. Choose the bottom most point in the middle box, this gives the
        mesh pattern
     \begin{equation*}
        \begin{tikzpicture}[scale=0.5]
            \modpattern[5]{}{1,2,3}{0/1,0/2,0/3,1/1,2/1,3/0,3/1,3/2}
            \draw (2.5,0.5) node {\(X\)};
        \end{tikzpicture}
    \end{equation*}
        Now we need to consider the box labelled \(X\). If this box is empty then
        we have an occurrence of \(m_2\) and are done. If this box contains
        any points then we gain some extra shading on the mesh pattern
        due to the dominating pattern
     \begin{equation*}
        \begin{tikzpicture}[scale=0.5]
            \modpattern[5]{2,3}{1,2,3}{0/1,0/2,0/3,1/1,1/2,1/3,2/1,3/0,3/1,3/2}
            \filldraw (2.5,0.5) circle (4 pt);
        \end{tikzpicture}
    \end{equation*}
        The two highlighted points form an occurrence of \(m_1\) with \(k\)
        points in the middle box, and thus by the Inductive Hypothesis we
        are done.
\end{description}
By induction we have that every occurrence of \(m_1\) leads to an occurrence of \(m_2\)
and by \NoteRef{not:downcmesh} every occurrence of \(m_2\) is an occurrence of
\(m_1\) so the two patterns are coincident.
This argument applies to another two pairs of classes. Therefore in
total in \(\av{\perm{3,2,1}, \mperm{1,2}{R}}\) there are \(213\) coincidence
classes.
